\scanpage{62}
Posons déjà
\[
(14)\qquad\left\{\begin{aligned}
V_{\Delta_1}&=\coprod_{\substack{(i,j)\in I^2\\i\leq j}}V_{i,j}\\
\overset{\deleted{\unsure{V}}}{V}_{\Delta_1}^*&=\coprod_{\substack{(i,j)\in I^2\\i\leq j}}V_{i,j}^*\underset{\text{ouvert}}{\subset}V_{\Delta_1}
\end{aligned}\right.
\]
on a donc
\diagram{(15)\qquad\begin{tikzcd}[column sep=large,row sep=large]
V_{\Delta_1}^*\arrow[r,hook]\arrow[d] & V_{\Delta_1}\arrow[d,dashed] & X_{\Delta_0}^*\arrow[l]\arrow[dl,dashed]\\
\unsure{X_{\Delta_0}^*|V_{\Delta_0}^*}\arrow[r,dashed] & X &
\arrow[from=2-1,to=2-2,phantom,"("{pos=.12}]\arrow[from=2-2,to=2-3,phantom,")"{pos=.05}]
\end{tikzcd}}
\diagram{\deleted{\unsure{(16)}\qquad\begin{tikzcd}
V_{\Delta_1} & X_{\Delta_0}^*\arrow[l]
\end{tikzcd}}}
qui reflètent les diagrammes
\diagram{(17)\qquad\begin{tikzpicture}
\node[inner sep=5mm] (d) {\begin{tikzcd}[column sep=large,row sep=large]
V_{i,j}^*\arrow[r]\arrow[d] & V_{ij}\arrow[d,dashed] & X_i^*\arrow[l]\arrow[dl,dashed]\\
X_j^*\arrow[r,dashed] & X &
\arrow[from=2-1,to=2-2,phantom,"("{pos=.12}]\arrow[from=2-2,to=2-3,phantom,")"{pos=.05}]
\end{tikzcd}};
\draw (d.north west) rectangle (d.south east);
\draw[->] (d.north west)--(d.west);
\end{tikzpicture}}

\margin{\uline{NB} La forme
\diagram{\begin{tikzcd}[column sep=small,row sep=small]
&\bullet\arrow[d]\\
\bullet\arrow[r]\arrow[d]&\bullet\\
\bullet&
\end{tikzcd}}
est peut être plus suggestive.}

\deleted{\unsure{On a}} Finalement, on n'est pas sûr que les notations (14) \unsure{soient|sont} bien utiles — de \unsure{tte|toute} façon le carré de (15) n'a rien de cocartésien, car le diagramme en trait plein, qui est le diagramme ``somme''
\mspage{61}
